\section{Theory}

\subsection{Deep Neural Networks}

Deep Neural Networks (DNNs) consist of an input layer, multiple hidden layers, and an output layer. Each 
layer contains nodes (neurons) connected to the next layer by weights $\theta$ and biases, which are 
updated during training to minimize a loss function.

\subsubsection{Universal Function Approximators}

A feedforward network with at least one hidden layer and a non-linear activation function can approximate 
any continuous function to arbitrary precision, given sufficient capacity (width). 

\subsubsection{Hierarchical Feature Representation}

As data passes through a deep network, each successive layer extracts increasingly abstract features. 
This essentially can be considered as automatic feature engineering. (=extracting critical components 
required for decision-making) It facilitates compression and therefore generalization.

\subsubsection{Scaling and Performance}

Increasing a networks width or its depth, increases its capacity to represent complex functions. While deeper 
networks are generally more parameter-efficient than wide, shallow ones, they are also harder to train. 

\subsection{Reinforcement Learning}

Reinforcement Learning is a type of Machine Learning where the agent learns by trial and error 
to optimize its behaviour for a return heuristic, as opposed to being trained on labeled examples. 
\\
\\
The return heuristic is calculated over episodes, which are sequences of interactions between agent and 
given environment, subject to evaluation. In RL episodes is modelled by Markov Decision Processes. 
We call specific realizations of episodes trajectories.

\subsubsection{Markov Decision Process - MDP}

The MDP is a sequence of states $s_i$, actions $a_i$ and rewards $r_i$

\[
s_1, a_1, r_1, s_2, a_2, r_2, s_3, a_3, r_3, \dots
\]

where each action $a_i$ is taken (in some way) based on the state $s_i$, results in a state transition 
$s_i \rightarrow s_{i+1}$ (controlled by the environment) and the reward $r_i$ (comes from environment).

In particular the MDP must follow Markov's Property

\paragraph{Markov's Property}
Observations are markovian if each subsequent state solely depends on the previous state and an action 
or in other words, the previous state and taken action are sufficient to determine the subsequent state. 
\[
P(s_{t+1} | s_t, a_t) = P(s_{t+1} | s_t, a_t, s_{t-1}, a_{t-1}, \dots, s_0, a_0)
\]

\subsubsection{Return $G_t$}

The return at $(s_t, a_t)$ is a function of the subsequent rewards: 

\[
    G_t = G_t(r_t, r_{t+1}, \dots)
\]

For example it can be defined as discounted sum of future rewards $G_t = \sum_{k=0}^{\infty} \gamma^k r_{t+k+1}$
, with $\gamma \in [0, 1]$.

\subsubsection{Policy $\pi(a_i|s_i)$}

The models objective is to learn of a policy $\pi(a_i|s_i)$, which for each $(s_t, a_t)$ picks an 
action such the respective expected return is maximized. (Environment may be stochastic)

\subsubsection{World Models}

The environment's internal mechanics are defined by the transition dynamics $P(s_{t+1}|s_t, a_t)$ 
and the reward function $R(s_t, a_t)$. Together, these are known as the World Model.

\begin{itemize}
    \item \textbf{Model-Free RL:} No access to transition and reward functions. Agent learns strictly from experience.
    \item \textbf{Model-Based RL:} Agent learns or is given a World Model. Thus it is able to simulate future trajectories without interacting with the environment (=planning).
\end{itemize}

\subsubsection{Value Functions}

To solve an MDP, the agent must estimate the "goodness" of being in a particular state or taking 
a particular action. This can be formalized with value functions, which estimate the expected return.

\paragraph{State-Value Function $V^\pi(s)$}

\begin{equation}
V^\pi(s) = \mathbb{E}_\pi [G_t | s_t = s]
\end{equation}

\paragraph{Action-Value Function (Q-Function) $Q^\pi(s, a)$}

\begin{equation}
Q^\pi(s, a) = \mathbb{E}_\pi [G_t | s_t = s, a_t = a]
\end{equation}

\paragraph{Optimal Value Functions $V^*$ and $Q^*$}
\begin{equation}
Q^*(s, a) = \max_\pi Q^\pi(s, a)
\end{equation}
and $V^*$ analogously.
\\
For Value-Based methods like DQN, the Q-function is the primary focus as it allows the agent 
to rank actions without needing a model of the environment's transitions. 
Optimal Value Functions resemble the "goal" of RL- They represent the maximum expected 
return achievable by any policy: If an agent learns $Q^*$, it can derive the optimal policy 
by simply selecting the action that maximizes the Q-value in any given state: $\pi^*(s) = \arg\max_a Q^*(s, a)$.

\subsubsection{Bellman Equation}
Value functions can be decomposed into a recursive (=Bellman Equation). It expresses the 
relationship between the value of a state (or state-action pair) and the value of its successor states.
For $Q^*$, the Bellman optimality equation is given by:
\begin{equation}
    Q^*(s, a) = \mathbb{E} [r_{t+1} + \gamma \max_{a'} Q^*(s_{t+1}, a') | s_t=s, a_t=a]
\end{equation}
This equation is the foundation of the DQN algorithm, as it provides a way to iteratively 
update the estimate of $Q^*$ based on immediate rewards and the predicted future value.

\subsubsection{Learning Methods}
There are two primary ways an agent can learn to estimate the value functions from 
experience: Monte Carlo methods and Temporal Difference learning.

\paragraph{Monte Carlo (MC) Methods}
Monte Carlo methods learn from complete episodes. The agent interacts with the environment 
until a terminal state is reached, calculates the total return $G_t$, and then updates the value 
of the states visited. While MC is unbiased, it requires the task to be episodic and can suffer 
from high variance.

\paragraph{Temporal Difference (TD) Learning}
TD learning combines ideas from Monte Carlo and Dynamic Programming. Unlike MC, TD methods 
can learn from incomplete sequences by "bootstrapping" a guess based on another guess. 

In TD(0), the simplest form of TD learning, the update target for a state $s$ is the immediate 
reward plus the discounted value of the next state:
\begin{equation}
    Target = r_{t+1} + \gamma V(s_{t+1})
\end{equation}

DQN utilizes TD learning, as it updates its Q-network at every time step using 
the transition $(s, a, r, s')$, making it significantly more sample-efficient than Monte Carlo approaches.

\subsection{Deep Q-Networks (DQN)}

Deep Q-Networks represent the fusion of the above. On the comparative to non-neural (tabular) RL, DQN win, 
because they are able to compress information (through generalization), whereas tabular RL would have to 
store all transitions in a LUT. This equivalently means that DQNs are suitable for continous state spaces, 
since they can estimate $Q(s, a; \theta) \approx Q^*(s, a)$, allowing to act correctly on previously unseen 
information. Tabular RL would simply fail due to lack of lookup values.
\\
\\
One major pitfall of DQN (but also tabular RL), is that despite accepting continuous state spaces, their outputs 
are discrete. In applications like robotics this can lead to jitter and lack of fine grained porecision.

\subsubsection{DQN Objective and Training}
DQN is a Value-Based, Model-Free algorithm that utilizes Temporal Difference learning. The network is trained by minimizing a loss function derived from the Bellman optimality equation. At each iteration $i$, the Mean Squared Error (MSE) of the TD error is minimized:
\begin{equation}
    L_i(\theta_i) = \mathbb{E} \left[ \left( \underbrace{r + \gamma \max_{a'} Q(s', a'; \theta_i^-)}_{\text{TD Target}} - Q(s, a; \theta_i) \right)^2 \right]
\end{equation}
Where $\theta_i^-$ represents the parameters of a "Target Network," a periodic snapshot of the weights used to stabilize training.

\subsubsection{Action Selection and Exploration}
Since the agent in DQN learns from its own experiences, it faces the \textit{exploration-exploitation} dilemma. Occasionally it is reasonable, to take sub-optimal actions to discover better long-term strategies, in order to learn action-value function.

\paragraph{Greedy Policy}
A strictly greedy policy always selects the action with the highest estimated value: $a = \arg\max_{a'} Q(s, a')$. This maximizes immediate reward, but prevents the agent from discovering better actions (lack of exploration).

\paragraph{Epsilon-Greedy ($\epsilon$-greedy)}
In this standard approach, the agent selects a random action with probability $\epsilon$ and the greedy action with probability $1-\epsilon$. Often, $\epsilon$ is annealed (decreased) over time, moving from high exploration to high exploitation as the agent's confidence in its Q-estimates grows.

\paragraph{Boltzmann (Softmax) Exploration}
Instead of choosing purely at random, Boltzmann exploration assigns a probability to each action based on its Q-value using a temperature parameter $\tau$:
\begin{equation}
    P(a|s) = \frac{e^{Q(s, a)/\tau}}{\sum_{a'} e^{Q(s, a')/\tau}}
\end{equation}
This ensures that actions with higher estimated values are explored more frequently than those with very low values.

\subsubsection{Experience Replay Buffer}
A core challenge in using DNNs for RL is that sequential samples in a trajectory are highly correlated, violating the \textit{i.i.d.} assumption required for stable gradient descent. Furthermore, the agent might "forget" rare but important experiences once it moves to a different part of the state space.

To address this, DQN utilizes an \textbf{Experience Replay Buffer} $\mathcal{D}$. 
\begin{enumerate}
    \item During interaction, agent stores transitions $e_t = (s_t, a_t, r_{t+1}, s_{t+1})$ in $\mathcal{D}$.
    \item During training, instead of picking most recent transition, agent samples a "mini-batch" of experiences uniformly at random from the buffer.
\end{enumerate}

This mechanism breaks the temporal correlation of data. It allows the network to learn from each experience multiple times, significantly improving sample efficiency and numerical stability.

\subsubsection{Target Networks and the Moving Target Problem}
A significant reason for instability in Deep Q-Learning arises from the same network parameters $\theta$ beeing used to calculate both the current estimate $Q(s, a; \theta)$ and the TD target. Because the target depends on the network's own predictions:
\begin{equation}
    Target = r + \gamma \max_{a'} Q(s', a'; \theta)
\end{equation}
any update to $\theta$ that reduces the error for the current state also changes the target for the next state. This creates a feedback loop where the "ground truth" the network is trying to learn is constantly shifting, leading to oscillations or divergence.

To solve this \textit{moving target problem}, DQN introduces a second neural network called the \textbf{Target Network}, with its own parameters $\theta^-$. 

\paragraph{Mechanism}
The target network is an identical copy of the online (policy) network but its weights are "frozen." Its exclusive use is to calculate the TD target:
\begin{equation}
    y^{DQN} = r + \gamma \max_{a'} Q(s', a'; \theta^-)
\end{equation}
The online network $\theta$ is updated at every training step, while the target network $\theta^-$ remains stationary. Periodically (every $C$ steps), the parameters of the online network are copied to the target network: $\theta^- \leftarrow \theta$.

\paragraph{Advantage}
By keeping the target fixed for a substantial number of iterations, the reinforcement learning problem is temporarily reduced to a supervised learning task with a stable objective. This means that wieght updates are decoupled from target calculation, which reduces oscillations and instability.